\subsection{serial}\label{serial}
The \verb|serial| module encapsulates serial input/output over UART
(optionally USB-CDC). It serves two purposes: first, it redirects the standard
output -- the newlib hooks \verb|_write| and \verb|_read| bind \verb|printf| to
the UART or USB so that debug output goes out directly. Second, it is itself a
\verb|kybd_t| implementation (Section \ref{keyboard}, device \verb|serial_dev|)
that turns received characters into key input within a \verb|state_t|.

The module connects several other building blocks: \verb|buffer| and
\verb|buffer_pool| (Sections \ref{buffer}, \ref{buffer_pool}) for the RX/TX
buffers, \verb|_time| (Section \ref{time}) for throughput measurement
(\verb|srxhdl|, \verb|stxhdl|) and \verb|cycle| (Section \ref{cycle}) for
timestamps.

\subsubsection*{Operating modes}
The behaviour is controlled via the bit field \verb|print_e| (Listing
\ref{print_e}) and set with \verb|serial_mode_set()|. It selects, among other
things, the transport layer (\verb|USE_UART|, \verb|USE_DMA_RX|,
\verb|USE_USB|) and optional preparation such as timestamps (\verb|TIMESTAMP|),
gap detection (\verb|GAP_DETECT|) or plain raw output (\verb|RAW|).

\begin{listing}
\begin{minted}{C}
typedef enum {
    RAW        = 1,
    TIMESTAMP  = 2,
    GAP_DETECT = 4,
    ONE_SHOT   = 0x8,
    USE_UART   = 0x10,
    USE_DMA_RX = 0x20,
    USE_DMA_TX = 0x40, // does not work
    USE_USB    = 0x80,
    MEASURE_BYTE_PER_SECONDS = 0x100,
} print_e;
\end{minted}
\caption{Output operating modes \texttt{print\_e}}
\label{print_e}
\end{listing}

\subsubsection*{State}
The configuration is passed at initialisation as \verb|sio_t| (Listing
\ref{sio_t}): UART handle, RX/TX buffers, a \verb|cycle_t| and the operating
mode. The module keeps the run-time state internally in the single variable
\verb|isio| (type \verb|isio_t|); this additionally holds the device state
\verb|state_t|, a \verb|buffer_pool_t| and counters for overflows and dropped
USB bytes.

\begin{listing}
\begin{minted}{C}
typedef struct _sio_t {
    UART_HandleTypeDef *uart;
    buffer_t *buffer[SIO_RXTX_CNT]; // [SIO_RX], [SIO_TX]
    cycle_t  *cycle;
    print_e   mode;
} sio_t;
\end{minted}
\caption{Configuration structure \texttt{sio\_t}}
\label{sio_t}
\end{listing}

\subsubsection*{Reception}
{\sloppy\emergencystretch=3em
Reception runs via DMA with idle-line detection
(\verb|HAL_UARTEx_ReceiveToIdle_DMA|). On reception the HAL calls the
\verb|HAL_UARTEx_RxEventCallback|: this copies the data into the RX buffer,
converts it to lowercase (\verb|buffer_tolower|) and passes it to the
\verb|state_t| via \verb|serial_apply_change()| -- an input recognised as a hex
number toggles the corresponding key.
\par}

\subsubsection*{API}
{\sloppy\emergencystretch=3em
\begin{description}
\item[\texttt{serial\_mode\_set(mode)}] Sets the \verb|print_e| bit field.
  Returns: \verb|void|.
\item[\texttt{serial\_mode\_get()}] Reads the \verb|print_e| bit field.
  Returns: \verb|print_e|.
\item[\texttt{serial\_reset(dev)}] Resets the RX/TX buffers and the device
  state. Returns: \verb|void|.
\item[\texttt{serial\_waitForNumber(key)}] Evaluates the last received input;
  \verb|key| then points to the string. Returns: \verb|int8_t| -- the number,
  or $-1$.
\item[\texttt{serial\_get\_byte\_per\_second()}] Reads the measured transmit
  throughput. Returns: \verb|uint32_t|.
\item[\texttt{serial\_reset\_byte\_per\_second()}] Resets the measured transmit
  throughput. Returns: \verb|void|.
\item[\texttt{serial\_uart\_printf(line, len)}] Prints \verb|len| bytes
  directly over the UART. Returns: \verb|uint32_t| (number of bytes sent).
\end{description}
\par}

\noindent Through the exported \verb|kybd_t| \verb|serial_dev| the module can
at the same time be used as an input device on the \verb|keyboard| interface;
the structure \verb|device_t serial_io| additionally offers generic
open/read/write access.
